\begin{frame}[label=caesarex, fragile]{Wiederholung letztes Mal: Angriff auf Caesar}
	\centering
	Wie würden Sie folgenden Text knacken?\\[.5cm]
\small\begin{verbatim}
BUYOZCAMNYLVOWBMNUVYCHXYONMWBYHNYRNYHCMNGYCMNYHMY...
\end{verbatim}
	    
	 \begin{itemize}[<+->]
	   \item Häufigster Buchstabe \textcolor{RoyalBlue}{im deutschen Alphabet}: \texttt{E} (= Buchstabe 4)
	   \item Häufigster Buchstabe im obigen \textcolor{ForestGreen}{im obigen Geheimtext}: \texttt{K} (= Buchstabe 10)
	   \item Schlüssel $=$ Verschiebung von $10-4 = 6 =$ Buchstabe \texttt{G}
	 \end{itemize}	 
		
	 \foreach \i in {1,...,7} {
	   \pgfmathtruncatemacro{\j}{\i + 3}
	   \pgfmathtruncatemacro{\shiftCaesar}{\i - 1}
	   \only<\j>{\centering\scalebox{.7}{\caesar[\shiftCaesar]}}
	 }
	 \only<10->{Können Sie den Text nun knacken?}
\end{frame}

\begin{frame}[fragile]{}
\centering
	Wie würden Sie folgenden Text knacken?\\[.5cm]
\begin{verbatim}
  Klartext: JEMANDMUSSTEJOSEFKVERLEUMDETHA...
 Schlüssel: GGGGGGGGGGGGGGGGGGGGGGGGGGGGGG...
Geheimtext: PKSGTJSAYYZKPUYKLQBKXRKASJKZNG...
\end{verbatim}
\end{frame}
